\documentclass[11pt]{article}
\usepackage[a4paper,margin=1.08in]{geometry}
\usepackage[T1]{fontenc}
\usepackage[utf8]{inputenc}
\usepackage{lmodern}
\usepackage{amsmath,amssymb,amsthm,mathtools}
\usepackage{enumitem}
\usepackage[hidelinks]{hyperref}
\usepackage{microtype}

\title{Foundations of Presentation Complexity}
\author{Luca Blanchi}
\date{June 2026}

\begin{document}
\maketitle

\begin{abstract}
This note lays out a foundational formalism for presentation complexity. It treats mathematical objects through chosen description systems, resource scales, bounded parts, cost-controlled transfers, simulations between presentation systems, and observables used for lower-bound and reconstruction arguments. The framework separates object presentation cost from observation cost and is intended as a common language for height theory, Kolmogorov complexity, circuit complexity, group presentations, algebraic straight-line complexity, and related AI-assisted notes in this collection.
\end{abstract}

\section{Guiding idea}

Many mathematical objects come with natural ways of being presented: by coordinates, generators and relations, equations, circuits, words, modules, kernels, local models, certificates, or finite approximations.

A presentation-complexity theory studies three related questions.

First, which descriptions are allowed, and how much do they cost?

Second, how does this cost behave under natural operations?

Third, what can be observed from objects that admit low-cost descriptions?

The basic principle is:

\[
\begin{gathered}
\text{cheap descriptions produce constrained objects,}\\
\text{and constrained objects produce constrained observations.}
\end{gathered}
\]

Conversely, if an object has an observation incompatible with every low-cost object, then the object cannot have a low-cost presentation.

The theory does not assign an absolute complexity to mathematical objects. It assigns complexity relative to a specified presentation system.

A second guiding principle is that presentation and observation should not be conflated. A description produces or specifies an object. An observable extracts information from an object. Many arguments use observables to study presentation complexity, but observables are not themselves presentation costs unless they arise from a presentation system.


\section{Object classes}

An \textbf{object class} is a class

\[
\mathcal X
\]

of mathematical objects.

Depending on the application, \(\mathcal X\) may be a set of objects, a set of isomorphism classes, or a category/groupoid whose objects are studied up to an equivalence relation.

Examples include:

\[
\mathcal X=\mathbb P^n(\mathbb Q),
\]

\[
\mathcal X={\text{finite groups up to isomorphism}},
\]

\[
\mathcal X={\text{Boolean functions on }n\text{ variables}},
\]

\[
\mathcal X={\text{polynomials over a fixed field}}.
\]

The theory begins only after choosing what counts as the object class.

This choice is not cosmetic. A finite group considered as an abstract group, a marked group, a permutation group, a matrix group, or a group equipped with structural data may lead to different presentation systems and different costs.


\section{Resource scales}

A \textbf{resource scale} is a preordered set

\[
(B,\leq)
\]

whose elements represent possible resource bounds.

The simplest case is

\[
B=\mathbb N
\]

with the usual order.

More generally, (B) may be a multi-resource scale such as

\[
B=\mathbb N^k
\]

with coordinatewise order.

This allows one to record several resources at once, such as size, depth, height, degree, arity, number of generators, number of relations, or total relation length.

A \textbf{cost bound} is an element

\[
b\in B.
\]

Informally, (b) is the available budget.

The theory should not force every cost into a single number. In many applications the bounded part is defined by several simultaneous constraints.


\section{Presentation systems}

A \textbf{presentation system} for an object class \(\mathcal X\) consists of a triple

\[
\Gamma=(\mathcal D,\rho,\kappa),
\]

where:

\[
\mathcal D
\]

is a class of descriptions,

\[
\rho:\mathcal D\to \mathcal X
\]

is a realization map, and

\[
\kappa:\mathcal D\to B
\]

is a cost function.

An element

\[
d\in\mathcal D
\]

is a description. The object realized by (d) is

\[
\rho(d).
\]

The cost of the description is

\[
\kappa(d).
\]

The presentation system tells us which descriptions are allowed and how expensive they are.

The same object class may admit many different presentation systems. This plurality is central to the theory.

For example, finite groups may be presented by generators and relations, multiplication tables, permutation actions, matrix representations, semidirect-product data, module data, commutator tensors, normal covers, or finite local models.

Different presentation systems measure different mathematical resources.


\section{Bounded descriptions and bounded objects}

For a cost bound \(b\in B\), define the bounded description class

\[
\mathcal D^{\leq b}=\{d\in\mathcal D:\kappa(d)\leq b\}.
\]

The corresponding bounded part of \(\mathcal X\) is

\[
\mathcal X_\Gamma^{\leq b}=\rho(\mathcal D^{\leq b}).
\]

Thus

\[
\mathcal X_\Gamma^{\leq b}
\]

is the class of objects admitting at least one \(\Gamma\)-description of cost at most (b).

Equivalently, an object \(x\in\mathcal X\) is (b)-presentable if

\[
x\in \mathcal X_\Gamma^{\leq b}.
\]

If \(b\leq b'\), then

\[
\mathcal X_\Gamma^{\leq b}\subseteq \mathcal X_\Gamma^{\leq b'}.
\]

Hence the bounded parts form a filtration of \(\mathcal X\):

\[
\mathcal X_\Gamma^{\leq b}
\subseteq
\mathcal X_\Gamma^{\leq b'}
\quad
\text{whenever }b\leq b'.
\]

This is the \textbf{presentation-complexity filtration} associated with \(\Gamma\).

The filtration is often the primary object. A single numerical complexity value may be less informative than the entire family of bounded parts.


\section{Presentation complexity}

When the resource scale admits a meaningful notion of infimum or minimum, one may define the \textbf{presentation complexity} of an object \(x\in\mathcal X\) by

\[
\operatorname{pc}_\Gamma(x)=\inf\{\kappa(d): d\in\mathcal D,\ \rho(d)=x\}.
\]

In discrete settings this is often a minimum:

\[
\operatorname{pc}_\Gamma(x)=\min\{\kappa(d):\rho(d)=x\}.
\]

Thus \(\operatorname{pc}_\Gamma(x)\) is the cost of the cheapest allowed description of \(x\).

The bounded part can then be written as

\[
\mathcal X_\Gamma^{\leq b}=\{x\in\mathcal X:\operatorname{pc}_\Gamma(x)\leq b\},
\]

whenever the minimum is attained or the inequality is interpreted as the existence of a description of cost at most (b).

The primary object is often the filtration

\[
b\mapsto \mathcal X_\Gamma^{\leq b},
\]

rather than the numerical value \(\operatorname{pc}_\Gamma(x)\) alone.


\section{Basic upper bound principle}

If \(d\in\mathcal D\) is a description of (x), meaning

\[
\rho(d)=x,
\]

then

\[
\operatorname{pc}_\Gamma(x)\leq \kappa(d).
\]

Thus every explicit description gives an upper bound on presentation complexity.

This is the constructive half of the theory.


\section{Cost-controlled maps}

Let

\[
\Gamma_\mathcal X=(\mathcal D_\mathcal X,\rho_\mathcal X,\kappa_\mathcal X)
\]

be a presentation system for \(\mathcal X\), and let

\[
\Gamma_\mathcal Y=(\mathcal D_\mathcal Y,\rho_\mathcal Y,\kappa_\mathcal Y)
\]

be a presentation system for \(\mathcal Y\).

Let

\[
F:\mathcal X\to\mathcal Y
\]

be a map of object classes.

A \textbf{description-level transfer} for (F) is a map

\[
\widehat F:\mathcal D_\mathcal X\to\mathcal D_\mathcal Y
\]

such that

\[
\rho_\mathcal Y(\widehat F(d))=F(\rho_\mathcal X(d))
\]

for every \(d\in\mathcal D_\mathcal X\).

The transfer is \textbf{cost-controlled with overhead \(f:B_\mathcal X\to B_\mathcal Y\)} if

\[
\kappa_\mathcal Y(\widehat F(d))
\leq
f(\kappa_\mathcal X(d))
\]

for every \(d\in\mathcal D_\mathcal X\).

In that case,

\[
F(\mathcal X_{\Gamma_\mathcal X}^{\leq b})
\subseteq
\mathcal Y_{\Gamma_\mathcal Y}^{\leq f(b)}.
\]

When presentation complexities are defined, this gives

\[
\operatorname{pc}_{\Gamma_\mathcal Y}(F(x))
\leq
f(\operatorname{pc}_{\Gamma_\mathcal X}(x)).
\]

Thus cost-controlled transfers formalize the idea that natural operations transport low-cost descriptions to low-cost descriptions.


\section{Cost-controlled operations}

More generally, let

\[
F:\mathcal X_1\times\cdots\times \mathcal X_n\to \mathcal Y
\]

be an (n)-ary operation.

A description-level transfer for (F) is a map

\[
\widehat F:
\mathcal D_1\times\cdots\times \mathcal D_n
\to
\mathcal D_\mathcal Y
\]

such that

\[
\rho_\mathcal Y(\widehat F(d_1,\dots,d_n))=F(\rho_1(d_1),\dots,\rho_n(d_n)).
\]

It is cost-controlled with overhead (f) if

\[
\kappa_\mathcal Y(\widehat F(d_1,\dots,d_n))
\leq
f(\kappa_1(d_1),\dots,\kappa_n(d_n)).
\]

Then

\[
\operatorname{pc}_{\Gamma_\mathcal Y}(F(x_1,\dots,x_n))
\leq
f(\operatorname{pc}_{\Gamma_1}(x_1),\dots,\operatorname{pc}_{\Gamma_n}(x_n))
\]

whenever the corresponding presentation complexities are defined.

This is the basic form of a \textbf{cost calculus}.

Examples include products, quotients, extensions, tensor products, composition of maps, base change, blow-ups, reductions modulo primes, and passage to associated graded or local models.


\section{Simulation between presentation systems}

Let

\[
\Gamma=(\mathcal D_\Gamma,\rho_\Gamma,\kappa_\Gamma)
\]

and

\[
\Delta=(\mathcal D_\Delta,\rho_\Delta,\kappa_\Delta)
\]

be presentation systems for the same object class \(\mathcal X\).

A \textbf{translation} from \(\Gamma\) to \(\Delta\) with overhead (f) is a map

\[
T:\mathcal D_\Gamma\to\mathcal D_\Delta
\]

such that

\[
\rho_\Delta(T(d))=\rho_\Gamma(d)
\]

and

\[
\kappa_\Delta(T(d))\leq f(\kappa_\Gamma(d)).
\]

Then

\[
\mathcal X_\Gamma^{\leq b}
\subseteq
\mathcal X_\Delta^{\leq f(b)}.
\]

When numerical presentation complexity is defined,

\[
\operatorname{pc}_\Delta(x)
\leq
f(\operatorname{pc}_\Gamma(x)).
\]

Thus \(\Delta\) simulates \(\Gamma\) up to overhead (f).

If there are translations in both directions with controlled overheads, then the two presentation systems are equivalent up to the corresponding resource distortion.

Separations between presentation systems arise when a family (\(x_n\)) has low complexity in one system and high complexity in another.


\section{Counting and volume}

The bounded parts define a counting function whenever cardinality is meaningful:

\[
b\mapsto |\mathcal X_\Gamma^{\leq b}|.
\]

This is the \textbf{volume profile} of the presentation system.

Counting arguments give generic lower bounds. If a family \(\mathcal S\subseteq\mathcal X\) satisfies

\[
|\mathcal S|>|\mathcal D^{\leq b}|,
\]

then some \(x\in\mathcal S\) has no description of cost at most (b).

This follows from

\[
|\mathcal X_\Gamma^{\leq b}|
\leq
|\mathcal D^{\leq b}|.
\]

Similarly, if

\[
|\omega(\mathcal S)|>|\operatorname{Val}_{\Gamma,\omega}(b)|,
\]

then some \(x\in\mathcal S\) has no description of cost at most (b).

Thus volume and observable volume provide nonconstructive lower bounds.

Counting arguments are useful for proving the existence of complex objects, but they usually do not identify a specific complex object unless combined with additional structure.


\section{Observables}

Let \(\mathcal X\) be an object class.

An \textbf{observable} on \(\mathcal X\) is a map

\[
\omega:\mathcal X\to A
\]

to an observation space (A).

The value

\[
\omega(x)
\]

is the information about (x) seen by the observable.

The observation space (A) need not be numerical. It may consist of numbers, vectors, modules, lattices, distributions, ranks, graphs, isomorphism classes, hom-counts, reductions modulo primes, local data, cohomology groups, or other structured data.

An observable is not the same thing as a cost. A cost measures the expense of a description. An observable measures information extracted from the object.

A numerical invariant may be an observable without being a presentation complexity.


\section{Attainable observations below a bound}

Let

\[
\Gamma=(\mathcal D,\rho,\kappa)
\]

be a presentation system for \(\mathcal X\), and let

\[
\omega:\mathcal X\to A
\]

be an observable.

For a cost bound \(b\in B\), define

\[
\operatorname{Val}_{\Gamma,\omega}(b)=\omega(\mathcal X_\Gamma^{\leq b})=\{\omega(x):x\in\mathcal X_\Gamma^{\leq b}\}.
\]

This is the set of \textbf{attainable observations below (b)}.

Equivalently, it is the set of observable values attained by objects admitting a \(\Gamma\)-description of cost at most (b).

If \(b\leq b'\), then

\[
\operatorname{Val}_{\Gamma,\omega}(b)
\subseteq
\operatorname{Val}_{\Gamma,\omega}(b').
\]

Thus

\[
b\mapsto \operatorname{Val}_{\Gamma,\omega}(b)
\]

is an increasing profile of observable values.

The older informal term "shadow" should be replaced by:

\[
\text{attainable observations below }b.
\]


\section{Observable profiles}

The \textbf{observable profile} of \(\omega\) relative to \(\Gamma\) is the function

\[
b\longmapsto \operatorname{Val}_{\Gamma,\omega}(b).
\]

It records how the observable sees the presentation-complexity filtration.

Instead of the full set \(\operatorname{Val}_{\Gamma,\omega}(b)\), one may study numerical summaries, such as

\[
b\mapsto |\operatorname{Val}_{\Gamma,\omega}(b)|,
\]

\[
b\mapsto \dim \operatorname{span}(\operatorname{Val}_{\Gamma,\omega}(b)),
\]

or other sizes, ranks, volumes, supports, or distributional statistics.

However, the full observable profile is the primary object.


\section{Observable lower-bound principle}

Let

\[
x\in\mathcal X.
\]

If

\[
\omega(x)\notin \operatorname{Val}_{\Gamma,\omega}(b),
\]

then

\[
x\notin \mathcal X_\Gamma^{\leq b}.
\]

Equivalently, (x) has no \(\Gamma\)-description of cost at most (b).

When presentation complexity is numerical, this gives

\[
\operatorname{pc}_\Gamma(x)> b.
\]

This is the basic observable lower-bound principle.


\section{Cost-observable bounds}

In practice, one often cannot determine the full set

\[
\operatorname{Val}_{\Gamma,\omega}(b).
\]

Instead, one proves an inclusion

\[
\operatorname{Val}_{\Gamma,\omega}(b)\subseteq S_b
\]

for some explicit subset

\[
S_b\subseteq A.
\]

Then

\[
\omega(x)\notin S_b
\]

implies

\[
x\notin \mathcal X_\Gamma^{\leq b}.
\]

Thus

\[
\omega(x)\notin S_b
\quad\Longrightarrow\quad
\operatorname{pc}_\Gamma(x)> b
\]

whenever numerical presentation complexity is defined.

This is the usual form of a lower-bound certificate:

\[
\text{low-cost descriptions force observations in }S_b;
\]

\[
\text{the object has observation outside }S_b;
\]

therefore the object is not low-cost.


\section{Observed sizes}

Suppose the observation space (A) has a size function

\[
s:A\to R
\]

to another ordered resource scale (R). Then

\[
s\circ\omega:\mathcal X\to R
\]

is an \textbf{observed size}.

It is a numerical or resource-valued invariant of objects.

However, an observed size is not automatically a presentation complexity.

It becomes a presentation complexity only if it is induced by some presentation system.

The important mathematical question is often to relate observed sizes to presentation complexity. For example, if one proves

\[
x\in\mathcal X_\Gamma^{\leq b}
\quad\Longrightarrow\quad
s(\omega(x))\leq F(b),
\]

then

\[
s(\omega(x))>F(b)
\quad\Longrightarrow\quad
x\notin\mathcal X_\Gamma^{\leq b}.
\]

When numerical presentation complexity is defined,

\[
s(\omega(x))>F(b)
\quad\Longrightarrow\quad
\operatorname{pc}_\Gamma(x)> b.
\]

This covers many familiar lower-bound arguments: degree too large, rank too large, denominator too large, dimension too large, relation lattice too constrained, or observable distribution outside the attainable range.

A useful rule is:

\[
\text{do not call an invariant a presentation cost unless it is induced by a presentation system.}
\]

But an invariant may still be a useful observed size.


\section{Families of observables}

A \textbf{family of observables} on \(\mathcal X\) is a collection

\[
\Omega={\omega_i:\mathcal X\to A_i}_{i\in I}.
\]

It may be viewed as a joint observable

\[
\omega_\Omega:\mathcal X\to \prod_{i\in I}A_i
\]

defined by

\[
\omega_\Omega(x)=(\omega_i(x))_{i\in I}.
\]

The family \(\Omega\) is \textbf{complete on a subclass \(\mathcal C\subseteq\mathcal X\)} if, for all \(x,y\in\mathcal C\),

\[
\omega_i(x)=\omega_i(y)\quad\text{for every }i\in I
\]

implies

\[
x=y,
\]

or, in an isomorphism-based context,

\[
x\cong y.
\]

A central problem is to determine which families of observables are complete on which bounded parts or natural subclasses.


\section{Indistinguishability fibres}

Given an observable

\[
\omega:\mathcal X\to A,
\]

the fibre

\[
\omega^{-1}(a)
\]

is the class of objects indistinguishable by \(\omega\) from the observable value (a).

If

\[
\omega(x)=\omega(y),
\]

then any method depending only on \(\omega\) gives the same output on (x) and (y).

Thus every observable has an intrinsic limitation: it can only distinguish objects up to its fibres.

For a family of observables \(\Omega\), define

\[
x\sim_\Omega y
\]

if

\[
\omega(x)=\omega(y)
\]

for every \(\omega\in\Omega\).

The equivalence classes of \(\sim_\Omega\) are the \textbf{indistinguishability fibres} of \(\Omega\).

Large fibres indicate blindness; small or trivial fibres indicate completeness.

Studying these fibres is essential when observables are used for classification.


\section{Reconstruction by observables}

A family of observables may do more than distinguish objects. It may reconstruct a structural invariant.

Let

\[
\tau:\mathcal X\to T
\]

be a structural invariant, such as a module, tensor, rank profile, support profile, lattice, isomorphism class, or moduli parameter.

A family \(\Omega\) \textbf{reconstructs} \(\tau\) on \(\mathcal C\subseteq\mathcal X\) if

\[
x\sim_\Omega y
\quad\Longrightarrow\quad
\tau(x)=\tau(y)
\]

for all \(x,y\in\mathcal C\).

Equivalently, \(\tau\) factors through the observational quotient

\[
\mathcal C/!\sim_\Omega.
\]

If \(\tau\) is a complete invariant on \(\mathcal C\), then reconstruction of \(\tau\) implies classification.

Many concrete results have this form:

\[
\text{word distributions reconstruct a rank profile,}
\]

\[
\text{commutator moments reconstruct a tensor pencil,}
\]

\[
\text{hom-count profiles reconstruct a module,}
\]

\[
\text{finite differences reconstruct coefficient data.}
\]

Thus reconstruction is a central use of observables beyond lower bounds.


\section{Costed observables}

Sometimes not only objects but also observations have costs.

A \textbf{costed observable family} on \(\mathcal X\) is a triple

\[
\mathfrak O=(\Omega,\chi,R),
\]

where:

\[
\Omega
\]

is a class of observables,

\[
R
\]

is an observation-resource scale, and

\[
\chi:\Omega\to R
\]

assigns an observation cost to each observable.

For an observation bound \(r\in R\), define

\[
\Omega^{\leq r}=\{\omega\in\Omega:\chi(\omega)\leq r\}.
\]

Thus \(\Omega^{\leq r}\) is the family of observables available with observation budget (r).

Examples of observation costs include arity of a word observer, length of a word, number or order of test groups, degree of a polynomial probe, rank or matrix size of a linear observer, number of evaluations, depth of a logical formula, or quantifier rank.

This distinction is essential:

\[
\text{presentation cost of objects}
\]

and

\[
\text{observation cost of tests}
\]

are different resources.

Costed observables should not replace object presentation costs. They form a second layer of the theory, needed when the act of observing, testing, or distinguishing also has complexity.


\section{Distinguishing cost}

Let

\[
\mathfrak O=(\Omega,\chi,R)
\]

be a costed observable family.

For two objects \(x,y\in\mathcal X\), define their \textbf{distinguishing cost} by

\[
\operatorname{dc}_{\mathfrak O}(x,y)=\inf\{\chi(\omega):\omega\in\Omega,\ \omega(x)\neq\omega(y)\}.
\]

If no observable in \(\Omega\) distinguishes (x) and (y), then

\[
\operatorname{dc}_{\mathfrak O}(x,y)=\infty.
\]

This measures the minimum observation budget needed to tell (x) and (y) apart using the allowed observables.

Observer-barrier results have the form:

\[
\Omega^{\leq r}\text{ is not complete on }\mathcal C,
\]

but

\[
\Omega^{\leq r'}\text{ is complete on }\mathcal C
\]

for some larger (r'), often with (r') optimal.

Such results show that observation itself has irreducible resource requirements.


\section{Normal-form compilers}

Many applications depend on a theorem that converts raw descriptions or raw observables into a simpler algebraic normal form.

A \textbf{normal-form compiler} is a result showing that a class of descriptions or observables is equivalent to a more structured class.

For descriptions, a normal-form compiler may have the form:

\[
d\in\mathcal D^{\leq b}
\quad\Longrightarrow\quad
\exists d'\in\mathcal D_{\mathrm{nf}}^{\leq f(b)}
\text{ with }
\rho(d')=\rho(d).
\]

This says that every low-cost description can be replaced by a normal-form description with controlled overhead.

For observables, a normal-form compiler may have the form

\[
\omega_{\mathrm{raw}}
\rightsquigarrow
\omega_{\mathrm{nf}},
\]

where the raw observable and the normal-form observable determine the same information on the class under study.

Examples include:

\[
\text{word equations in class-two groups}
\rightsquigarrow
\text{linear and alternating quadratic data,}
\]

\[
\text{relative word observers}
\rightsquigarrow
\text{matrix-rank observers,}
\]

\[
\text{Boolean lift constraints}
\rightsquigarrow
\text{Möbius coefficient constraints,}
\]

\[
\text{commutator-word moments}
\rightsquigarrow
\text{rank probes of alternating tensors,}
\]

\[
\text{coordinate descriptions}
\rightsquigarrow
\text{primitive height-normalized coordinates.}
\]

Normal-form compilers are often the technical heart of applications. They make bounded objects and attainable observations computable, constrainable, or reconstructive.


\section{Realization images and algebraic obstructions}

A presentation system contains a realization map

\[
\rho:\mathcal D\to\mathcal X.
\]

Often the key problem is to understand the image of \(\rho\), or the image of a bounded, evaluated, linearized, or reduced version of \(\rho\).

For a bound (b), the bounded image is

\[
\mathcal X_\Gamma^{\leq b}=\rho(\mathcal D^{\leq b}).
\]

In algebraic settings, obstructions may arise from failure to lie in an image, cokernel classes, congruence conditions, rank constraints, lattice constraints, support constraints, or coefficient alphabet restrictions.

Thus some lower bounds are not initially about size. They are about realizability.

The general pattern is:

\[
\text{low-cost descriptions land inside a constrained image;}
\]

\[
\text{the target object or observation lies outside that image;}
\]

therefore

\[
\text{no low-cost description exists.}
\]

This is especially important when the presentation system is algebraic, linear, modular, or arithmetic.


\section{Pullback and transport of observables}

Let

\[
F:\mathcal X\to\mathcal Y
\]

be a map of object classes, and let

\[
\eta:\mathcal Y\to A
\]

be an observable on \(\mathcal Y\).

Then

\[
\omega=\eta\circ F:\mathcal X\to A
\]

is an observable on \(\mathcal X\).

This is the \textbf{pullback observable} induced by (F) and \(\eta\).

If (F) is cost-controlled with overhead (f), then

\[
F(\mathcal X_{\Gamma_\mathcal X}^{\leq b})
\subseteq
\mathcal Y_{\Gamma_\mathcal Y}^{\leq f(b)}.
\]

Therefore,

\[
\operatorname{Val}_{\Gamma_\mathcal X,\eta\circ F}(b)
\subseteq
\operatorname{Val}_{\Gamma_\mathcal Y,\eta}(f(b)).
\]

Thus maps between object classes transport observable profiles.

This is useful when (F(x)) is a quotient, reduction, abelianization, associated graded object, local model, or simpler invariant of (x).

A transported lower bound has the following form.

If

\[
\operatorname{pc}_{\Gamma_\mathcal Y}(F(x))> f(b),
\]

then

\[
\operatorname{pc}_{\Gamma_\mathcal X}(x)> b.
\]

Thus a lower bound for a simpler transported object can imply a lower bound for the original object.


\section{Relative and fibre methods}

Sometimes an object is studied through a map

\[
q:\mathcal X\to\mathcal Y.
\]

The object (q(x)) may be a quotient, reduction, base object, abelianization, associated graded object, invariant, or observable component of (x).

The fibre

\[
q^{-1}(q(x))
\]

contains the information not determined by (q(x)).

One may try to decompose the cost of (x) into the cost of the base object (q(x)) and the residual cost of specifying (x) inside the fibre.

Informally,

\[
\text{cost of }x
\approx
\text{cost of }q(x)
+
\text{cost of lifting data}.
\]

This method is useful for extensions, quotients, fibrations, local-to-global problems, model-plus-residual descriptions, cohomological classification, and situations where an observable sees only part of the object.

A full relative theory may introduce conditional presentation systems or relative costs

\[
\operatorname{pc}_\Gamma(x\mid q(x)).
\]

But the core idea is already present in the basic formalism through fibres of maps and observables.


\section{Bounded-counterexample arguments}

Let (P) be a property of objects in \(\mathcal X\). Define the bad set

\[
B_P={x\in\mathcal X : P(x)\text{ fails}}.
\]

Thus proving (P) for all objects in \(\mathcal X\) is the same as proving

\[
B_P=\varnothing.
\]

Let \(\Gamma\) be a presentation system on \(\mathcal X\), and let (b) be a cost bound.

A \textbf{bounded-counterexample principle} for (P), with respect to \(\Gamma\) and (b), is an implication of the form

\[
B_P\neq\varnothing
\quad\Longrightarrow\quad
B_P\cap \mathcal X_\Gamma^{\leq b}\neq\varnothing.
\]

In words: if (P) fails somewhere, then (P) already fails on an object admitting a \(\Gamma\)-presentation of cost at most (b).

This principle alone does not prove (P). It only says that any failure of (P) can be witnessed by a low-cost object.

To prove (P), one combines it with an exclusion result:

\[
B_P\cap \mathcal X_\Gamma^{\leq b}=\varnothing.
\]

Together, the two statements imply

\[
B_P=\varnothing.
\]

Indeed, if \(B_P\neq\varnothing\), the bounded-counterexample principle gives an element of

\[
B_P\cap \mathcal X_\Gamma^{\leq b},
\]

contradicting the exclusion result.

The exclusion result may be proved directly, by finite enumeration, by a normal-form theorem, or by observables.

The observable form is the following. Suppose there is an observable

\[
\omega:\mathcal X\to A
\]

and a subset \(S_b\subseteq A\) such that all objects of cost at most (b) have observable value in \(S_b\):

\[
\mathcal X_\Gamma^{\leq b}\subseteq \omega^{-1}(S_b),
\]

while every bad object has observable value outside \(S_b\):

\[
B_P\subseteq \mathcal X\setminus \omega^{-1}(S_b).
\]

Then no bad object can lie in the bounded part. If bad objects, if they exist, must lie in the bounded part, then no bad objects exist.

In slogan form:

\[
\text{counterexamples must be cheap}
+
\text{counterexamples cannot be cheap}
\Rightarrow
\text{no counterexamples}.
\]

This argument should not be confused with the false heuristic that high-complexity objects cannot exist. High-complexity objects may certainly exist. The point is different: if a counterexample, by an independent mathematical reason, would have to admit a low-cost presentation, but low-cost objects are excluded from being counterexamples, then no counterexample exists.


\section{Admissibility and naturality}

The formal definition of a presentation system is intentionally broad. Not every formal presentation system is mathematically meaningful.

A presentation system should normally be required to satisfy some admissibility conditions appropriate to the application.

Common conditions include:

These are not part of the minimal formalism, but they are necessary for the theory to avoid artificial examples.

A central methodological rule is:

\[
\text{do not call an invariant a cost unless it is induced by a presentation system.}
\]

Likewise:

\[
\begin{gathered}
\text{do not use an observable for lower bounds unless its choice is itself admissible}\\
\text{or its cost is accounted for.}
\end{gathered}
\]


\section{Basic examples}


\subsection{Heights}

For rational points in projective space,

\[
\mathcal X=\mathbb P^n(\mathbb Q),
\]

one may take descriptions to be primitive integer coordinate vectors

\[
(a_0,\dots,a_n)\in \mathbb Z^{n+1}_{\mathrm{prim}}\setminus{0},
\]

with realization

\[
\rho(a_0,\dots,a_n)=[a_0:\dots:a_n]
\]

and cost

\[
\kappa(a_0,\dots,a_n)=\log\max_i|a_i|.
\]

The resulting presentation complexity is the logarithmic height.

Thus height theory is an arithmetic instance of presentation complexity.

Height balls are bounded parts. Height inequalities under morphisms are cost-controlled maps. Counting points of bounded height is volume theory for bounded parts.


\subsection{Kolmogorov complexity}

For finite binary strings,

\[
\mathcal X={0,1}^*,
\]

take descriptions to be programs for a universal machine (U):

\[
\mathcal D={0,1}^*,
\]

with realization

\[
\rho(p)=U(p)
\]

and cost

\[
\kappa(p)=|p|.
\]

Then

\[
K_U(x)=\min\{|p|:U(p)=x\}
\]

is Kolmogorov complexity relative to (U).


\subsection{Circuit complexity}

For Boolean functions,

\[
\mathcal X={f:{0,1}^n\to{0,1}},
\]

descriptions may be Boolean circuits, realization is the computed function, and cost is circuit size or depth.

The presentation complexity is minimum circuit size or depth.


\subsection{Group presentations}

For finite groups, one possible presentation system uses finite presentations

\[
\langle S\mid R\rangle,
\]

with realization

\[
\rho(S,R)=F(S)/\langle!\langle R\rangle!\rangle
\]

and cost such as

\[
\kappa(S,R)=\bigl(|S|,|R|,\sum_{r\in R}|r|\bigr).
\]

Another presentation system may use multiplication tables, word circuits, semidirect-product data, module data, commutator tensors, finite local models, or permutation-cover data.

Different presentation systems measure different mathematical resources.

Words may play two distinct roles. As relators, words are part of descriptions and contribute to presentation cost. As word maps or word equations, words define observables and contribute to observation cost.


\subsection{Algebraic and straight-line complexity}

For polynomials or algebraic expressions, descriptions may be arithmetic circuits or straight-line programs.

The realization map sends a circuit to the polynomial or function it computes.

The cost may be the number of operations, depth, degree of allowed gates, or circuit size.

Observables may include degree, rank of partial derivative spaces, support, evaluation data, or factorization patterns.

This gives a presentation-complexity view of algebraic complexity.


\section{What the theory proves}

The basic theory gives a collection of reusable proof schemes.


\subsection*{Construction}

To prove

\[
\operatorname{pc}_\Gamma(x)\leq b,
\]

construct a description (d) such that

\[
\rho(d)=x
\]

and

\[
\kappa(d)\leq b.
\]


\subsection*{Cost propagation}

To prove that an operation preserves low complexity, construct a description-level transfer

\[
\widehat F
\]

and bound its cost overhead.


\subsection*{Simulation}

To compare two presentation systems, construct translations between their description spaces and bound the overhead.


\subsection*{Separation}

To separate two presentation systems \(\Gamma\) and \(\Delta\), find objects \(x_n\) such that

\[
\operatorname{pc}_\Gamma(x_n)
\]

is small while

\[
\operatorname{pc}_\Delta(x_n)
\]

is large.


\subsection*{Observable obstruction}

To prove

\[
\operatorname{pc}_\Gamma(x)> b,
\]

find an observable \(\omega\) and a constraint \(S_b\) such that

\[
\operatorname{Val}_{\Gamma,\omega}(b)\subseteq S_b
\]

but

\[
\omega(x)\notin S_b.
\]


\subsection*{Counting}

To prove the existence of high-complexity objects, compare the number of low-cost descriptions with the number of objects or observable values under consideration.


\subsection*{Classification}

To classify a bounded part or natural subclass, find a family of observables complete on that class.


\subsection*{Blindness}

To prove that an observable family is insufficient, construct non-isomorphic objects that the family cannot distinguish.


\subsection*{Reconstruction}

To explain what an observable family sees, show that it reconstructs a structural invariant.


\subsection*{Observer-cost barrier}

To prove that observation itself has irreducible complexity, show that observables below a certain observation budget are not complete, while observables above that budget are complete.


\subsection*{Normal form}

To make bounded objects or observations analyzable, prove that arbitrary descriptions or observables can be compiled into a controlled normal form.


\subsection*{Realization obstruction}

To prove that an object is not low-cost, show that low-cost descriptions land in a constrained image, and that the target object or its observation lies outside that image.


\subsection*{Transported lower bound}

To lower-bound (x), lower-bound a cost-controlled image (F(x)).


\subsection*{Bounded-counterexample argument}

To prove a universal statement (P), show that any counterexample would have to lie in a bounded part, and then exclude all bad objects from that bounded part.


\subsection*{Relative or fibre argument}

To understand the cost of (x), decompose it into the cost of a base object (q(x)) and the residual cost of lifting within the fibre over (q(x)).


\section{Scope}

The minimal theory contains:

\[
\text{presentation systems},
\]

\[
\text{bounded parts},
\]

\[
\text{presentation-complexity filtrations},
\]

\[
\text{cost-controlled transfers},
\]

\[
\text{simulations between presentation systems},
\]

\[
\text{volume profiles},
\]

\[
\text{observables},
\]

\[
\text{attainable observations},
\]

\[
\text{observable profiles},
\]

\[
\text{lower-bound certificates},
\]

\[
\text{families of observables},
\]

\[
\text{indistinguishability fibres}.
\]

The following are not separate foundations, but important core methods and refinements:

\[
\text{observed sizes},
\]

\[
\text{costed observables},
\]

\[
\text{distinguishing cost},
\]

\[
\text{observable reconstruction},
\]

\[
\text{normal-form compilers},
\]

\[
\text{realization-image obstructions},
\]

\[
\text{relative or fibre methods}.
\]

Further modules may add approximate presentations, probabilistic or statistical certificates, dynamics of complexity under iteration, and more detailed genericity theories.

These modules should be introduced only when demanded by applications.

The guiding discipline is that abstraction should serve proofs: every new notion should support at least one of construction, cost control, lower bounds, classification, separation, observer limitation, reconstruction, or finite refutation.

The central hierarchy is:

\[
\text{object presentation systems first,}
\]

\[
\text{cost calculus second,}
\]

\[
\text{observables as tools third,}
\]

\[
\text{observation cost as a further layer when needed.}
\]

This keeps the theory centered on presentation complexity while still allowing the observer-based methods that drive many applications.

\section*{Declaration of generative AI and AI-assisted technologies in the writing process}
During the preparation of this work, ChatGPT, by OpenAI, was used to assist with mathematical drafting, formalization, review, and editing. This work is shared as a preliminary AI-assisted mathematical note. The mathematical content may have been only partially reviewed and may contain errors; it should not be treated as peer-reviewed or as a fully verified manuscript.

\begin{thebibliography}{9}
\bibitem{AroraBarak2009}
S. Arora and B. Barak,
\emph{Computational Complexity: A Modern Approach},
Cambridge University Press, 2009.

\bibitem{Kolmogorov1965}
A. N. Kolmogorov,
Three approaches to the quantitative definition of information,
\emph{Problems of Information Transmission} 1 (1965), 1--7.

\bibitem{LiVitanyi2008}
M. Li and P. Vitanyi,
\emph{An Introduction to Kolmogorov Complexity and Its Applications},
third edition, Springer, 2008.

\bibitem{Northcott1950}
D. G. Northcott,
An inequality in the theory of arithmetic on algebraic varieties,
\emph{Proceedings of the Cambridge Philosophical Society} 45 (1950), 502--509.
\end{thebibliography}

\end{document}
